% Options for packages loaded elsewhere
\PassOptionsToPackage{unicode}{hyperref}
\PassOptionsToPackage{hyphens}{url}
%
\documentclass[
  11pt,
  letterpaper,
]{article}
\usepackage{amsmath,amssymb}
\usepackage{iftex}
\ifPDFTeX
  \usepackage[T1]{fontenc}
  \usepackage[utf8]{inputenc}
  \usepackage{textcomp} % provide euro and other symbols
\else % if luatex or xetex
  \usepackage{unicode-math} % this also loads fontspec
  \defaultfontfeatures{Scale=MatchLowercase}
  \defaultfontfeatures[\rmfamily]{Ligatures=TeX,Scale=1}
\fi
\usepackage{lmodern}
\ifPDFTeX\else
  % xetex/luatex font selection
\fi
% Use upquote if available, for straight quotes in verbatim environments
\IfFileExists{upquote.sty}{\usepackage{upquote}}{}
\IfFileExists{microtype.sty}{% use microtype if available
  \usepackage[]{microtype}
  \UseMicrotypeSet[protrusion]{basicmath} % disable protrusion for tt fonts
}{}
\makeatletter
\@ifundefined{KOMAClassName}{% if non-KOMA class
  \IfFileExists{parskip.sty}{%
    \usepackage{parskip}
  }{% else
    \setlength{\parindent}{0pt}
    \setlength{\parskip}{6pt plus 2pt minus 1pt}}
}{% if KOMA class
  \KOMAoptions{parskip=half}}
\makeatother
\usepackage{xcolor}
\usepackage[top=2.5cm,bottom=2.5cm,left=2.5cm,right=2.5cm]{geometry}
\usepackage{longtable,booktabs,array}
\usepackage{calc} % for calculating minipage widths
% Correct order of tables after \paragraph or \subparagraph
\usepackage{etoolbox}
\makeatletter
\patchcmd\longtable{\par}{\if@noskipsec\mbox{}\fi\par}{}{}
\makeatother
% Allow footnotes in longtable head/foot
\IfFileExists{footnotehyper.sty}{\usepackage{footnotehyper}}{\usepackage{footnote}}
\makesavenoteenv{longtable}
\usepackage{graphicx}
\makeatletter
\def\maxwidth{\ifdim\Gin@nat@width>\linewidth\linewidth\else\Gin@nat@width\fi}
\def\maxheight{\ifdim\Gin@nat@height>\textheight\textheight\else\Gin@nat@height\fi}
\makeatother
% Scale images if necessary, so that they will not overflow the page
% margins by default, and it is still possible to overwrite the defaults
% using explicit options in \includegraphics[width, height, ...]{}
\setkeys{Gin}{width=\maxwidth,height=\maxheight,keepaspectratio}
% Set default figure placement to htbp
\makeatletter
\def\fps@figure{htbp}
\makeatother
\setlength{\emergencystretch}{3em} % prevent overfull lines
\providecommand{\tightlist}{%
  \setlength{\itemsep}{0pt}\setlength{\parskip}{0pt}}
\setcounter{secnumdepth}{-\maxdimen} % remove section numbering
\ifLuaTeX
\usepackage[bidi=basic]{babel}
\else
\usepackage[bidi=default]{babel}
\fi
\babelprovide[main,import]{spanish}
% get rid of language-specific shorthands (see #6817):
\let\LanguageShortHands\languageshorthands
\def\languageshorthands#1{}
% Soporte para caracteres Unicode (modo una columna)
\usepackage{newunicodechar}

% Definir caracteres Unicode problemáticos
\newunicodechar{₀}{\ensuremath{_0}}
\newunicodechar{₁}{\ensuremath{_1}}
\newunicodechar{₂}{\ensuremath{_2}}
\newunicodechar{₃}{\ensuremath{_3}}
\newunicodechar{₄}{\ensuremath{_4}}
\newunicodechar{₅}{\ensuremath{_5}}
\newunicodechar{₆}{\ensuremath{_6}}
\newunicodechar{₇}{\ensuremath{_7}}
\newunicodechar{₈}{\ensuremath{_8}}
\newunicodechar{₉}{\ensuremath{_9}}
\newunicodechar{⁺}{\ensuremath{^+}}
\newunicodechar{⁻}{\ensuremath{^-}}
\newunicodechar{°}{\textdegree}
\newunicodechar{×}{\texttimes}
\newunicodechar{·}{\textperiodcentered}
\newunicodechar{–}{\textendash}
\newunicodechar{—}{\textemdash}
\newunicodechar{λ}{\ensuremath{\lambda}}
\newunicodechar{σ}{\ensuremath{\sigma}}
\newunicodechar{ε}{\ensuremath{\varepsilon}}
\newunicodechar{μ}{\ensuremath{\mu}}
\newunicodechar{χ}{\ensuremath{\chi}}
\newunicodechar{α}{\ensuremath{\alpha}}
\newunicodechar{β}{\ensuremath{\beta}}
\newunicodechar{γ}{\ensuremath{\gamma}}
\newunicodechar{δ}{\ensuremath{\delta}}
\newunicodechar{θ}{\ensuremath{\theta}}
\newunicodechar{φ}{\ensuremath{\phi}}
\newunicodechar{ω}{\ensuremath{\omega}}
\newunicodechar{π}{\ensuremath{\pi}}
\newunicodechar{Σ}{\ensuremath{\Sigma}}
\newunicodechar{Δ}{\ensuremath{\Delta}}
\newunicodechar{∞}{\ensuremath{\infty}}
\newunicodechar{≈}{\ensuremath{\approx}}
\newunicodechar{≠}{\ensuremath{\neq}}
\newunicodechar{≤}{\ensuremath{\leq}}
\newunicodechar{≥}{\ensuremath{\geq}}
\newunicodechar{→}{\ensuremath{\rightarrow}}
\newunicodechar{←}{\ensuremath{\leftarrow}}
\newunicodechar{↔}{\ensuremath{\leftrightarrow}}
\newunicodechar{↑}{\ensuremath{\uparrow}}
\newunicodechar{↓}{\ensuremath{\downarrow}}
\newunicodechar{±}{\ensuremath{\pm}}
\newunicodechar{÷}{\ensuremath{\div}}
\newunicodechar{≪}{\ensuremath{\ll}}
\newunicodechar{≫}{\ensuremath{\gg}}
\newunicodechar{∇}{\ensuremath{\nabla}}
\newunicodechar{∂}{\ensuremath{\partial}}
\ifLuaTeX
  \usepackage{selnolig}  % disable illegal ligatures
\fi
\usepackage{bookmark}
\IfFileExists{xurl.sty}{\usepackage{xurl}}{} % add URL line breaks if available
\urlstyle{same}
\hypersetup{
  pdflang={es},
  hidelinks,
  pdfcreator={LaTeX via pandoc}}

\author{}
\date{}

\begin{document}

\section{De la Electrocultura a la Agricultura Plasma: Un Arco de Tres
Siglos que Conecta el Legado de Bertholon con los Avances Agrícolas
Contemporáneos}\label{de-la-electrocultura-a-la-agricultura-plasma-un-arco-de-tres-siglos-que-conecta-el-legado-de-bertholon-con-los-avances-agruxedcolas-contemporuxe1neos}

\textbf{Autor:} T. Dufour\\
\textbf{Afiliación:} Laboratoire de Physique des Plasmas, Sorbonne
Université, CNRS, École Polytechnique\\
\textbf{Revista:} Compte-Rendus Mécanique de l'Académie des Sciences,
Volumen 354 (2026), pp.~89-116\\
\textbf{DOI:} \url{https://doi.org/10.5802/crmeca.331}

\begin{center}\rule{0.5\linewidth}{0.5pt}\end{center}

\subsection{Resumen}\label{resumen}

Este artículo traza una línea continua desde las ``electrizaciones'' del
siglo XVIII hasta la agronomía de plasma frío de hoy. El eje es la
metrología: donde Bertholon diseñó caminos para un ``fluido eléctrico''
indivisible sin acceso a dosis ni agentes, la agricultura plasma
contemporánea define fuentes, cuantifica transferencias interfaciales y
mide la dosis bajo condiciones controladas.

\begin{center}\rule{0.5\linewidth}{0.5pt}\end{center}

\section{1. Electricidad y agronomía: una trayectoria histórica desde la
Antigüedad hasta el siglo
XXI}\label{electricidad-y-agronomuxeda-una-trayectoria-histuxf3rica-desde-la-antiguxfcedad-hasta-el-siglo-xxi}

\subsection{1.1. Electricidad (Antigüedad--siglo XVIII): del espectáculo
experimental a la instrumentación
metroológica}\label{electricidad-antiguxfcedadsiglo-xviii-del-espectuxe1culo-experimental-a-la-instrumentaciuxf3n-metrooluxf3gica}

Desde la Antigüedad hasta finales del siglo XVIII, el estudio de los
fenómenos eléctricos evolucionó de una curiosidad cualitativa a una
práctica cuantitativa basada en instrumentos. Autores tempranos como
Teofrasto (\emph{De lapidibus}) y Plinio el Viejo (\emph{Historia
Natural}) señalaron la atracción de la paja y otros objetos ligeros por
el ámbar frotado (\emph{elektron}) y el azabache (gagates), mientras que
autores médicos (p.~ej., Scribonius Largus, Galeno) reportaron el uso
analgésico de peces eléctricos (torpedos) \textbf{{[}1--3{]}}. Los
lapidarios y enciclopedias medievales transmitieron ampliamente estas
observaciones, siendo el \emph{De mineralibus} de Alberto Magno el
ejemplo emblemático \textbf{{[}4{]}}. Durante el Renacimiento, Girolamo
Cardano (\emph{De subtilitate}, 1550) y Giambattista della Porta
(\emph{Magia naturalis}, 1558/1589) ampliaron el inventario de ``cuerpos
eléctricos'' (ámbar, azabache, diamante, zafiro) capaces de atraer
pequeñas partículas tras la fricción \textbf{{[}5, 6{]}}.

La consolidación conceptual comenzó con el trabajo de William Gilbert
\emph{De Magnete} (1600), quien acuñó el término ``electricus'' y
distinguió claramente las atracciones eléctricas del magnetismo
\textbf{{[}7{]}}. En los siglos XVII y principios del XVIII, las
máquinas de fricción multiplicaron las demostraciones: el globo de
azufre de Otto von Guericke (década de 1660) y los globos de vidrio
luminiscentes de Francis Hauksbee (1705--1710) son ejemplos canónicos
\textbf{{[}8, 9{]}}. Los experimentos de Stephen Gray (1729)
establecieron la conducción y el aislamiento como comportamientos
distintos, y Charles-François de Cisternay du Fay (1733) organizó los
efectos bajo una dualidad ``vitrea/resinosa'' \textbf{{[}10, 11{]}}.

Solo en la década de 1740, los avances en instrumentación crearon
condiciones para experimentos acumulativos. La botella de Leyden,
desarrollada independientemente por E.-G. von Kleist y P. van
Musschenbroek (1745--1746), fue el primer dispositivo confiable para
almacenar carga eléctrica (Figura 1a). Conectar varias botellas en
paralelo para formar una ``batería'' aumentó aún más la capacidad de
almacenamiento (Figura 1b) \textbf{{[}12--14{]}}. Al mismo tiempo, las
máquinas electrostáticas se estandarizaron, ejemplificadas por el
cilindro de Nairne (1773; Figura 1c) y el generador monumental de van
Marum (1784) \textbf{{[}15, 16{]}}. Las demostraciones públicas
proporcionaron un puente entre espectáculo y experimento: el 14 de marzo
de 1746, el Abbé Nollet descargó famosamente una botella de Leyden a
través de una cadena de guardias del rey en Versalles y luego a través
de monjes cartujos \textbf{{[}17, 18{]}}.

Las herramientas metroológicas se multiplicaron: la demostración de
Franklin de la equivalencia entre rayos y electricidad y su promoción
del pararrayos (1750--1752) \textbf{{[}17, 21, 22{]}}. El electrómetro
de lectura de chispa de Lane (1767) permitió la cuantificación de la
descarga de la botella de Leyden, mientras que el cuadrante de Henley
(1770) introdujo una medida de repulsión de un solo índice
\textbf{{[}23--25{]}}. H.-B. de Saussure extendió los electrómetros a
estudios atmosféricos durante las décadas de 1760--1780
\textbf{{[}26{]}}. Este giro cuantitativo culminó con la balanza de
torsión de Coulomb (1785), que produjo la primera ley precisa de fuerzas
eléctricas \textbf{{[}27{]}}. Obras de síntesis como la \emph{History
and Present State of Electricity} de Priestley (1767/1769), el
\emph{Complete Treatise of Electricity} de Cavallo (1777) y el
\emph{Dell'elettricismo artificiale e naturale} de Beccaria (1753)
organizaron la literatura en expansión \textbf{{[}14, 28--30{]}}.

Habiendo establecido herramientas para almacenamiento, entrega y
medición, las primeras aplicaciones a sistemas vivos comenzaron a
surgir. En 1783, \emph{De l'électricité des végétaux} de Bertholon
introdujo el ``electro-végétomètre,'' un dispositivo diseñado para
capturar y dirigir la electricidad atmosférica hacia los cultivos
\textbf{{[}31{]}}. Stephen Demainbray en la década de 1740 electrificó
mirto y observó crecimiento acelerado, resultados confirmados por el
Abbé Nollet en semillas de mostaza y Jean Jallabert en bulbos de jacinto
y narciso. Más tarde, Edouard-François Nunebert comparó bulbos de
cebolla no electrificados y electrificados, notando germinación más
rápida y tallos más altos \textbf{{[}31{]}}.

\begin{quote}
\includegraphics{./images/image1.png}
\end{quote}

\begin{quote}
\textbf{\emph{Figura 1. (a) Botella de Leyden: un recipiente de vidrio
que actúa como dieléctrico, con un revestimiento interior de hojalata
conectado por una cadena al terminal superior; una hoja exterior
(conectada a tierra) forma la segunda armadura {[}19{]}. (b) Batería de
botellas de Leyden (``batería de Franklin'') en una caja de madera
{[}19{]}. (c) Máquina eléctrica de cilindro de Nairne (c.~1770s)
{[}20{]}.}}
\end{quote}

\subsection{1.2. 1790--2000: Del galvanismo a las descargas controladas
y sus
procesos}\label{del-galvanismo-a-las-descargas-controladas-y-sus-procesos}

La invención de la pila por Volta (1799--1800; Figura 2a) marcó un punto
decisivo: por primera vez, los experimentadores tenían acceso a una
fuente de corriente continua, controlable y reproducible
\textbf{{[}32{]}}. A lo largo del siglo XIX, esta base se reforzó con la
estandarización de instrumentos y unidades: galvanómetros y puentes de
precisión, celdas estándar y el sistema práctico voltio-ohmio-amperio
proporcionaron un marco metroológico calibrado.

Faraday formalizó las leyes de la electrólisis en la década de 1830
\textbf{{[}36{]}}. En biología, Burdon-Sanderson registró potenciales de
acción en \emph{Dionaea muscipula} (1873), mostrando que los movimientos
vegetales podían asociarse con señales eléctricas discretas
\textbf{{[}37{]}}. Una generación después, Bose extendió este programa a
la electrofisiología vegetal, desarrollando instrumentos y métodos para
detectar, transmitir y analizar señales de plantas (Figura 2b)
\textbf{{[}33{]}}.

\begin{quote}
\includegraphics{./images/image2.jpeg}
\end{quote}

\begin{quote}
\textbf{\emph{Figura 2. (a) Pila voltaica (1800): columna de discos
alternos de zinc-cobre separados por arandelas empapadas en salmuera
{[}20{]}. (b) Aparato de Bose para respuestas mecánicas y eléctricas
simultáneas {[}33{]}. (c) Tubo de Crookes (alto vacío) que demuestra
rayos catódicos {[}34{]}. (d) Primer ozonizador de Werner von Siemens
(1857) {[}35{]}.}}
\end{quote}

Desde la segunda mitad del siglo XIX, comenzó a surgir una verdadera
metrología de las descargas. Paschen formuló su ley (1889) que
relacionaba el voltaje de ruptura con la presencia y la distancia entre
electrodos \textbf{{[}38{]}}. La teoría de ionización de Townsend
(1897--1901) aclaró la multiplicación en avalancha de electrones
\textbf{{[}39, 40{]}}. Los tubos de Geissler y los experimentos de
Crookes con gases enrarecidos (Figura 2c) revelaron la diversidad de
fenómenos luminosos a presión reducida \textbf{{[}34, 41{]}}. Ya en
1857, Werner von Siemens describió la descarga de barrera dieléctrica
(DBD) como fuente práctica para generación de ozono (Figura 2d)
\textbf{{[}35{]}}. Para 1928, Langmuir introdujo el término plasma para
describir un gas ionizado cuasi-neutro \textbf{{[}42{]}}.

Sobre esta base metroológica, siguieron aplicaciones químicas
estructuradas. Los ozonizadores DBD de Siemens permitieron la
ozonización de agua a gran escala \textbf{{[}43, 44{]}}. El arco
Birkeland--Eyde (1903, Figura 3a) aprovechó plasmas de alta temperatura
para oxidar nitrógeno atmosférico, produciendo nitratos como precursores
de fertilizantes \textbf{{[}45{]}}. Las lámparas de vapor de mercurio a
baja presión sistematizaron el uso germicida de la radiación
ultravioleta (Figura 3b) \textbf{{[}46--48{]}}.

\begin{quote}

\includegraphics{./images/image3.jpeg}
\end{quote}

\begin{quote}
\textbf{\emph{Figura 3. (a) Disco de llama Birkeland--Eyde {[}45{]}. (b)
Conducto de prueba de aerosol germicida UV {[}48{]}.}}
\end{quote}

A principios del siglo XX, la electrocultura experimentó un renovado
interés. Los dispositivos iban desde las redes aéreas de cables cargados
de Lemström hasta las antenas ``atmosféricas'' de Christofleau
\textbf{{[}49, 50{]}}. Como se muestra en la Figura 4a, estas últimas
estaban diseñadas para canalizar la ``electricidad natural'' de diversas
fuentes hacia el suelo. La Figura 4b representa su aplicación a viñedos,
mientras que la Figura 4c ilustra informes de crecimiento mejorado de
frijoles. Sin embargo, estos resultados deben interpretarse con
precaución: los regímenes de fertilización, la heterogeneidad del suelo,
los efectos estacionales y las geometrías de los dispositivos estaban
mal controlados.

\begin{quote}
\includegraphics{./images/image4.jpeg}
\end{quote}

\begin{quote}
\textbf{\emph{Figura 4. (a) Diagrama esquemático de la antena
atmosférica de Christofleau {[}50{]}. (b) Aplicación a vides {[}50{]}.
(c) Frijoles cultivados con estiércol (izquierda) vs.~bajo la antena de
Christofleau (derecha) {[}50{]}.}}
\end{quote}

Esta corriente histórica de electrocultura no evolucionó directamente
hacia tecnologías basadas en plasma. Estas surgieron independientemente
a finales del siglo XX, impulsadas por avances en física del plasma.
Desde la década de 1970, los plasmas fríos comenzaron a investigarse
como agentes de esterilización a baja temperatura. Para la década de
1990, los esterilizadores de plasma gaseoso de peróxido de hidrógeno
entraron en la práctica sanitaria \textbf{{[}51, 52{]}}. A
mediados-finales de la década de 1990, surgieron fuentes atmosféricas
(DBDs, reactores glidarc, chorros de plasma) que operaban directamente
al aire libre \textbf{{[}53{]}}.

\subsection{1.3. El siglo XXI: el auge de la agricultura
plasma}\label{el-siglo-xxi-el-auge-de-la-agricultura-plasma}

La agricultura plasma designa un campo de investigación contemporáneo
que aplica fuentes de plasma frío a semillas, plantas, suelos, sistemas
de riego y cadenas agroalimentarias. El análisis bibliométrico (Figura
5a) muestra que el término ``agricultura plasma'' ganó tracción solo
después de mediados de la década de 2010.

Los plasmas fríos son descargas no térmicas cuya eficacia surge de la
coacción de cinco propiedades (Figura 5b):

\begin{itemize}
\tightlist
\item
  \textbf{Eléctrica:} campos transitorios, frentes de ionización y
  corrientes pulsadas
\item
  \textbf{Química:} especies reactivas de oxígeno y nitrógeno (RONS) ---
  de vida corta (OH, O, NO) y de vida larga (H₂O₂, O₃, NO₂⁻/NO₃⁻)
\item
  \textbf{Radiativa:} fotones UV y VUV que se acoplan a objetivos de
  ADN/proteínas
\item
  \textbf{Térmica:} el gas volumétrico permanece cerca del ambiental,
  aunque el calentamiento local (10--50 °C) puede ser diseñado
\item
  \textbf{Mecánico-fluida:} viento iónico y microchorros mejoran la
  transferencia de masa en límites plasma-líquido
\end{itemize}

\begin{quote}
\includegraphics{./images/image5.png}
\end{quote}

\begin{quote}
\textbf{\emph{Figura 5. (a) Publicaciones anuales que contienen
``agricultura plasma'' (Google Scholar, 1990--2024). (b) Diagrama de
semillas tratadas por CAP con sus 5 propiedades acopladas.}}
\end{quote}

Para 2025, la agricultura plasma se organiza en cuatro ejes de
aplicación:

\begin{itemize}
\tightlist
\item
  \textbf{Tratamiento de semillas:} modificación superficial inducida
  por plasma y RONS/UV de baja dosis para acelerar la germinación
  \textbf{{[}57, 58{]}}. La química de PAW se cuantifica rutinariamente
  mediante trazadores de dosis \textbf{{[}59{]}}.
\item
  \textbf{Tratamiento de plantas:} CAP, afterglows y PAW bajo control
  térmico estricto para provocar señalización en lugar de daño
  \textbf{{[}60--62{]}}.
\item
  \textbf{Tratamiento de medios:} agua de riego, suelo y aire --- CAP y
  PAW permiten desinfección suave y modulación de comunidades
  microbianas \textbf{{[}63, 64{]}}.
\item
  \textbf{Tratamiento agroalimentario:} CAP y PAW para sanitizar
  productos, empaques y equipos \textbf{{[}65--67{]}}.
\end{itemize}

En la agricultura plasma contemporánea, el progreso se define por un
marco basado en la física de dosimetría explícita: métricas de dosis
cuantificables y rastreables (concentraciones de RONS, fluencia UV,
dosis eléctrica, tiempo de exposición) se miden en el objetivo.

\begin{center}\rule{0.5\linewidth}{0.5pt}\end{center}

\section{2. Reevaluación de la contribución del Abbé Bertholon a la
electricidad agrícola (década de
1780)}\label{reevaluaciuxf3n-de-la-contribuciuxf3n-del-abbuxe9-bertholon-a-la-electricidad-agruxedcola-duxe9cada-de-1780}

\subsection{2.1. Breve biografía
científica}\label{breve-biografuxeda-cientuxedfica}

Pierre-Nicolas Bertholon de Saint-Lazare (1741--1800), sacerdote, físico
y pedagogo de la Ilustración, construyó su carrera en Montpellier, donde
enseñó física experimental \textbf{{[}68{]}}. Durante la década de 1780
se convirtió en una figura central de la electricidad experimental,
notable por su enfoque en aplicaciones más que en espectáculo. Combinó
instrucción pública con demostraciones cuidadosamente preparadas,
enmarcando su programa en torno a la electricidad meteorológica, la
electricidad animal y la electricidad vegetal \textbf{{[}69{]}}.

Bertholon diseñó el electro-végétomètre, promovió la ``electricidad
médica'' y abogó por pararrayos en Languedoc. Participó en el movimiento
aerostático temprano, lanzando un globo aerostático con Jean-Antoine
Chaptal en 1783 \textbf{{[}69, 71{]}}.

\subsection{2.2. Tesis de Bertholon: la electricidad como estímulo para
la
vegetación}\label{tesis-de-bertholon-la-electricidad-como-estuxedmulo-para-la-vegetaciuxf3n}

Bertholon avanzó una tesis notable: que la electricidad atmosférica era
un reservorio constante capaz de estimular toda la fisiología de las
plantas --- germinación, circulación de la savia, nutrición, secreción,
crecimiento, floración, fructificación y transpiración
\textbf{{[}31{]}}. Su tratado se estructuró en tres partes:

\begin{enumerate}
\def\labelenumi{\arabic{enumi}.}
\tightlist
\item
  Demostrar la existencia e influencia de la electricidad atmosférica
\item
  Catalogar sus efectos en las funciones vegetales
\item
  Operationalizar estos efectos mediante instrumentos dedicados,
  principalmente el electro-végétomètre
\end{enumerate}

Reunió tres tipos de evidencia: correlaciones cualitativas (tormentas
coincidiendo con germinación más rápida), observaciones instrumentadas
(electrómetros, cometas y electroscopios) y conjeturas mecanicistas
(electricidad impulsando la circulación de savia a través de ``afinidad
eléctrica'').

La fortaleza del trabajo de Bertholon radica en su visión unificadora.
Los límites también eran claros: no se intentó dosimetría, los roles de
diferentes agentes no se distinguieron y gran parte del razonamiento
permaneció analógico.

\subsection{2.3. Dispositivos y procedimientos de operación de
electrocultura de
Bertholon}\label{dispositivos-y-procedimientos-de-operaciuxf3n-de-electrocultura-de-bertholon}

\subsubsection{Dispositivos de ruta seca: el
electro-végétomètre}\label{dispositivos-de-ruta-seca-el-electro-vuxe9guxe9tomuxe8tre}

El instrumento insignia de Bertholon fue el electro-végétomètre --- un
colector-distribuidor pasivo de ``electricidad atmosférica'' sin
pantalla ni dosimetría. Se diseñaron dos modelos (Figura 6):

\begin{quote}
\includegraphics{./images/image6.jpeg}
\end{quote}

\begin{quote}
\textbf{\emph{Figura 6. Electró-végétomètre de Bertholon. (a) Modelo I:
colector de punto único y brazo articulado. (b) Modelo II: corona de
puntos, brazo giratorio extensible. Ambos incluían un dispositivo de
neutralización (Z) {[}31{]}.}}
\end{quote}

\textbf{Tabla 1. Comparación de las variantes del electro-vegetómetro de
Bertholon.}

\begin{longtable}[]{@{}
  >{\raggedright\arraybackslash}p{(\columnwidth - 4\tabcolsep) * \real{0.1404}}
  >{\raggedright\arraybackslash}p{(\columnwidth - 4\tabcolsep) * \real{0.3684}}
  >{\raggedright\arraybackslash}p{(\columnwidth - 4\tabcolsep) * \real{0.4912}}@{}}
\toprule\noalign{}
\begin{minipage}[b]{\linewidth}\raggedright
Característica
\end{minipage} & \begin{minipage}[b]{\linewidth}\raggedright
Modelo I (punto único, brazo articulado)
\end{minipage} & \begin{minipage}[b]{\linewidth}\raggedright
Modelo II (corona de múltiples puntos, brazo giratorio)
\end{minipage} \\
\midrule\noalign{}
\endhead
\bottomrule\noalign{}
\endlastfoot
\textbf{Propósito} & Captura pasiva → Conducción aislada → Difusión
suave & Igual \\
\textbf{Mástil/Base} & Mástil de madera: sección enterrada secada al
fuego, alquitranada, envuelta en carbón/cemento & Mástil con cilindro de
cabeza saturado en resina, tratado con alquitrán/brea \\
\textbf{Aislamiento superior} & Tubo de vidrio grueso con mastic
bituminoso (D); segundo aislante (H) & Interfaz rotativa aislada
mediante inserción de resina (D) y manguito de vidrio/mastic (H) \\
\textbf{Colector} & Varilla de punto único (E) con embudo de hojalata
(F) & Corona de múltiples puntos (E) para maximizar captura de carga \\
\textbf{Ruta de conducción} & Cadena conductora (G) desde E a través de
H & Palanca de hierro curvada (C) que lleva cadena (G) al brazo \\
\textbf{Buffer} & Disco de hierro (K) como pequeño
``condensador/regulador'' & Sin disco separado; regulación implícita \\
\textbf{Brazo horizontal} & Conductor articulado (K-L-M-N) en caballetes
aislantes (O, P) con cuerdas de seda & Conductor telescópico (L-M) con
conductor telescópico (Q) en caballete aislante (P) \\
\textbf{Difusor terminal} & Corona de puntas afiladas (N, R) que liberan
corona no disruptiva & Corona de puntas afiladas \\
\textbf{Movilidad} & Brazo se balancea en bisagras (L, Q) para cubrir un
sector & Rotación completa de 360°; extender/retraer para variar
alcance \\
\textbf{Neutralización} & Cadena entre K y tierra para drenar carga &
Igual \\
\textbf{Seguridad} & ``Excitador en forma de C'' (marco de cobre/hierro
con mango de vidrio + cadena de tierra) & Igual \\
\end{longtable}

Ambos modelos incluían un ``excitador en forma de C'': una herramienta
rudimentaria de puesta a tierra para seguridad del operador (recuadro de
la Figura 6b). Consistía en una varilla metálica curvada con extremos
redondeados, uno de los cuales llevaba una cadena arrastrándose en el
suelo.

Bertholon también imaginó un método de ``ruta seca'' para combatir
insectos taladrantes de la madera (Figura 7). Propuso descargar una
botella de Leyden directamente a través de la madera afectada
posicionando dos alambres de hierro en lados opuestos. Para huertos, los
árboles se conectaban con alambres de hierro desnudos para que una sola
descarga pudiera propagarse a través de tramos sucesivos.

\begin{quote}
\includegraphics{./images/image7.jpeg}
\end{quote}

\begin{quote}
\textbf{\emph{Figura 7. (a) Huerto de frutales donde los árboles están
interconectados por alambres de hierro formando una cuadrícula. Una
botella de Leyden se descarga a través de la red. (b) Primer plano: dos
alambres enmarcan una sección del tronco {[}31{]}.}}
\end{quote}

\subsubsection{Dispositivos de ruta húmeda: riego
electrificado}\label{dispositivos-de-ruta-huxfameda-riego-electrificado}

Bertholon también diseñó instrumentos basados en riego con ``agua
electrificada,'' impulsados por una máquina electrostática (Figura 8).

\begin{quote}
\includegraphics{./images/image8.jpeg}
\end{quote}

\begin{quote}
\textbf{\emph{Figura 8. (Izquierda) ``Riego eléctrico artificial'' para
árboles: operador en taburete aislante con cubeta y bomba de jeringa.
(Derecha) ``Riego eléctrico'' móvil para camas de hortalizas: carrito
aislado con resina y carrete de alambre {[}31{]}.}}
\end{quote}

\textbf{Tabla 2. Comparación de los dispositivos de riego electrificado
de Bertholon.}

\begin{longtable}[]{@{}
  >{\raggedright\arraybackslash}p{(\columnwidth - 4\tabcolsep) * \real{0.2051}}
  >{\raggedright\arraybackslash}p{(\columnwidth - 4\tabcolsep) * \real{0.2821}}
  >{\raggedright\arraybackslash}p{(\columnwidth - 4\tabcolsep) * \real{0.5128}}@{}}
\toprule\noalign{}
\begin{minipage}[b]{\linewidth}\raggedright
Característica
\end{minipage} & \begin{minipage}[b]{\linewidth}\raggedright
Riego árbol por árbol
\end{minipage} & \begin{minipage}[b]{\linewidth}\raggedright
Carrito móvil para camas de hortalizas
\end{minipage} \\
\midrule\noalign{}
\endhead
\bottomrule\noalign{}
\endlastfoot
\textbf{Aislamiento} & Operador en taburete aislante (A), tabla
soportada por pilares de vidrio & Operador en carrito aislante moldeado
con resina (A) \\
\textbf{Fuente de agua} & Cubeta (B) sobre plataforma del taburete &
Regaderas rellenadas por segundo jardinero \\
\textbf{Dispositivo de riego} & Bomba de jeringa (C) rocía sobre dosel
(G) & Operador vierte desde regadera mientras sostiene cordón vivo
(G) \\
\textbf{Conexión eléctrica} & Máquina electrostática vía cadena (E)
desde conductor (D) & Carrito lleva carrete (D) con eje aislado, alambre
(C) a máquina (H) \\
\textbf{Conducción al agua} & Placa de hojalata (F) bajo el pie asegura
que la cubeta se comunique con la carga & Alambre pasa por soportes
aislantes (E) con anillos de vidrio (F) \\
\textbf{Modo de electrificación} & El agua se carga; chispas/hormigueo
en el chorro & Regadera y operador al mismo potencial; ``lluvia
eléctrica'' \\
\textbf{Movilidad/flujos} & Dispositivo transportado árbol por árbol;
simple e inexpensive & Carrito avanza fila por fila; suministro continuo
de regaderas \\
\textbf{Seguridad} & Aislamiento asegura que la corriente fluya al agua;
leve hormigueo si conectado a tierra & Operador aislado en carrito;
``excitador en forma de C'' recomendado \\
\textbf{Uso agronómico} & Efectivo cuando se repite diariamente sobre
plántulas & Ligeramente más lento pero beneficios claros tras varios
días \\
\end{longtable}

\subsection{2.4. Limitaciones técnicas y
metodológicas}\label{limitaciones-tuxe9cnicas-y-metodoluxf3gicas}

A finales del siglo XVIII, el programa de Bertholon estaba limitado
sobre todo por la instrumentación. El electrómetro de chispa de Lane
(1767) y el cuadrante de Henley (c.~1770) proporcionaron solo
indicadores cualitativos toscos \textbf{{[}72, 73{]}}. La balanza de
torsión de Coulomb (1785), aunque estableció la primera ley cuantitativa
de la electrostática, posterior a los escritos de Bertholon y midió
fuerzas en lugar de deposición de energía \textbf{{[}74, 75{]}}.

El marco conceptual era igualmente incompleto. El ozono no se reconoció
como especie química distinta hasta 1840 \textbf{{[}76{]}}, mientras que
la radiación ultravioleta se descubrió solo en 1801 \textbf{{[}77,
78{]}}. Las especies reactivas de oxígeno y nitrógeno (RONS) y los
flujos electrohidrodinámicos eran completamente desconocidos. En
\emph{De l'Électricité des Végétaux} (pp.~399, 401), el propio Bertholon
reportó aigrettes luminosas visibles en la noche en los terminales de
múltiples puntos --- la física moderna las reconoce como descargas
corona: plasmas localizados que liberan iones, electrones, ozono y
óxidos de nitrógeno \textbf{{[}31, 79{]}}.

La práctica experimental estaba muy por debajo de los estándares
modernos. Bertholon dependió de relatos cualitativos sin replicación
sistemática, randomización ni grupos de control formales. Las
herramientas metodológicas requeridas para estabilizar la inferencia
solo aparecerían en el siglo XX \textbf{{[}80, 81{]}}.

En contraste, la agricultura plasma moderna define la ``dosis'' a través
de parámetros explícitos y cuantitativos: intensidad del campo eléctrico
(V·m⁻¹·s), densidad de carga superficial (C·m⁻²) y deposición de energía
específica (J·kg⁻¹ o J·m⁻²) \textbf{{[}54, 56, 82{]}}.

\begin{center}\rule{0.5\linewidth}{0.5pt}\end{center}

\section{3. Física moderna de plasmas fríos aplicada a la
agricultura}\label{fuxedsica-moderna-de-plasmas-fruxedos-aplicada-a-la-agricultura}

\subsection{3.1. Definición de plasma frío y propiedades
clave}\label{definiciuxf3n-de-plasma-fruxedo-y-propiedades-clave}

El plasma es el cuarto estado de la materia, caracterizado por
cuasi-neutralidad macroscópica (densidades aproximadamente equilibradas
de iones positivos y electrones) y dinámica colectiva dominada por
interacciones de Coulomb de largo alcance. Dos parámetros fundamentales
gobiernan su comportamiento electrostático:

\begin{itemize}
\tightlist
\item
  \textbf{Longitud de Debye} λ\_D = √(ε₀ k\_B T\_e / n\_e e²)
\item
  \textbf{Frecuencia de plasma de electrones} ω\_p = √(n\_e e² / ε₀
  m\_e)
\end{itemize}

Un medio se califica como plasma cuando λ\_D ≪ L (tamaño del sistema),
N\_D = (4/3)πn\_e λ\_D³ ≫ 1, y los efectos colectivos dominan sobre
colisiones binarias \textbf{{[}83, 84{]}}.

Un \textbf{plasma atmosférico frío (CAP)} es una descarga de no
equilibrio en la que la temperatura de los electrones T\_e (≈1--5 eV, es
decir, 11,600--58,000 K) excede ampliamente la temperatura del gas T\_g
(≈300--350 K). A presión atmosférica, se observan regímenes de streamer
(microdescargas canalizadas \textless10 ns) y descargas de barrera
dieléctrica (DBDs), que pueden ser filamentosas o difusas
\textbf{{[}84{]}}.

Los CAP acoplan múltiples procesos: el campo eléctrico impulsa el
transporte de carga mientras que los RONS incluyen especies de vida
corta (O(³P), OH•, ¹O₂, NO•) y de vida larga (O₃, H₂O₂, NO₂⁻/NO₃⁻)
\textbf{{[}87, 88{]}}. La cadena de energía: fuente → electrones →
reacciones elementales → especies reactivas → flujos interfaciales
\textbf{{[}89{]}}.

Las fuentes atmosféricas comunes incluyen \textbf{DBDs} (planares,
coaxiales) impulsadas sinusoidalmente o en modo pulsado (kHz--MHz)
\textbf{{[}90{]}} y \textbf{chorros de plasma a presión atmosférica
(APPJ)} que operan en helio o argón \textbf{{[}91{]}}.

La cualificación metroológica incluye: diagnósticos eléctricos (V(t),
I(t), figuras de Lissajous), imagen ICCD para visualización de
streamers, temperatura del gas por espectroscopia rotacional y análisis
fisicoquímico (H₂O₂, NO₂⁻, NO₃⁻, pH, conductividad) \textbf{{[}87,
92--96{]}}.

\subsection{3.2. Tecnologías y protocolos de
laboratorio}\label{tecnologuxedas-y-protocolos-de-laboratorio}

El CAP se implementa a través de tres rutas complementarias:

\begin{itemize}
\tightlist
\item
  \textbf{DDS (Exposición directa en superficie seca):} plasma aplicado
  directamente a tejidos sin agua libre intencional
\item
  \textbf{DWS (Exposición directa en superficie húmeda):} plasma
  aplicado con una película acuosa delgada en la interfaz
\item
  \textbf{PAW (Agua activada por plasma):} agua tratada ex situ por
  plasma, luego aplicada por remojo, aspersión, riego o enjuague
  postcosecha \textbf{{[}58{]}}
\end{itemize}

Rangos operativos indicativos: - \textbf{DBD:} 6--20 kVpp, 5--50 kHz,
gap 0.5--5 mm, exposición 0.5--60 s - \textbf{APPJ:} 2--10 kVpp, 10--200
kHz o RF, 0.5--5 L·min⁻¹ He/Ar, boquilla-tejido 5--30 mm

Para exposiciones directas, la dosis se expresa como energía específica
(J·m⁻² para hojas, J·kg⁻¹ para lotes de semillas), condiciones
operativas eléctricas, composición/flujo del gas, distancia
fuente-objetivo, velocidad de escaneo, tiempo de exposición y
temperatura del efluente \textbf{{[}98, 99{]}}.

Para PAW, la caracterización incluye energía eléctrica específica
(J·L⁻¹), estado químico (H₂O₂, NO₂⁻/NO₃⁻, pH, conductividad) y dosis
aplicada \textbf{{[}87, 100, 101{]}}.

\subsection{3.3. Efectos del plasma frío en las
semillas}\label{efectos-del-plasma-fruxedo-en-las-semillas}

El CAP ejerce tres influencias principales e interdependientes sobre las
semillas: (1) alteración de la fisico-química de la cubierta, (2)
reducción de la carga microbiana y (3) activación de respuestas
fisiológicas y moleculares tempranas.

\subsubsection{3.3.1. Modificación de la cubierta de la semilla y
humectabilidad}\label{modificaciuxf3n-de-la-cubierta-de-la-semilla-y-humectabilidad}

Los estudios de microscopía (AFM/SEM) muestran consistentemente que el
CAP altera la superficie externa de las semillas mediante el alisamiento
y reestructuración de la cubierta a través de grabado, ablación y
nano/micro-estructuración. En trigo y cebada, los tratamientos con
plasma Ar--O₂ produjeron superficies visiblemente fisuradas y más
rugosas que se correlacionaron con una mayor absorción de agua
\textbf{{[}102, 103{]}}.

\begin{quote}
\includegraphics{./images/image9.jpeg}
\end{quote}

\begin{quote}
\textbf{\emph{Figura 9. Imágenes SEM de superficies de semillas de trigo
tratadas durante 90 s con (a) control (sin plasma), (b) plasma Ar/O₂.
Barra de escala es 50 μm. Reproducido de Rahman et al., Scientific
Reports, 2018, CC BY 4.0 {[}102{]}.}}
\end{quote}

La química de la superficie también se transforma: el plasma introduce
grupos funcionales que contienen oxígeno, aumentando la energía
superficial \textbf{{[}107, 108{]}}. Una de las señales más claras es la
reducción dramática en el ángulo de contacto con el agua (WCA). En
lenteja, frijol y trigo, el plasma de RF de aire redujo el WCA de 127°,
98° y 115° a 20°, 53° y casi 0°, respectivamente (Figura 10)
\textbf{{[}109{]}}.

\begin{quote}
\includegraphics{./images/image10.jpeg}
\end{quote}

\begin{quote}
\textbf{\emph{Figura 10. Gotita de agua sobre una semilla de frijol: (a)
sin tratar, (b) tratada con plasma. Reproducido de Bormashenko et al.,
Scientific Reports, 2012, CC BY 4.0 {[}109{]}.}}
\end{quote}

\subsubsection{3.3.2. Descontaminación y control de
patógenos}\label{descontaminaciuxf3n-y-control-de-patuxf3genos}

Los tratamientos con CAP logran reducciones sustanciales en la carga
microbiana superficial para maíz, trigo, girasol y otras semillas
\textbf{{[}111{]}}. La Figura 11a demuestra la descontaminación fúngica
efectiva de semillas de girasol tras un tratamiento de plasma de 5
minutos. En garbanzo, 2--5 minutos de CAP redujeron la microbiota
natural en 1--2 unidades log \textbf{{[}112{]}}. En avellanas, el plasma
de SF₆ a baja presión logró hasta una reducción de 5-log de
\emph{Aspergillus parasiticus} \textbf{{[}113{]}}.

\begin{quote}
\includegraphics{./images/image11.jpeg}
\end{quote}

\begin{quote}
\textbf{\emph{Figura 11. (a) Semillas de girasol contaminadas por
Rhizopus: sin tratar vs.~tratadas con plasma durante 5 min (T. Dufour,
LPP). (b) Efecto del CAP sobre la germinación de Cicer arietinum y la
supervivencia de microorganismos {[}112{]}. (c) Reducción de A.
parasiticus en avellanas {[}113{]}.}}
\end{quote}

Sin embargo, la inactivación microbiana completa sigue siendo
infrecuente y dependiente de la especie. En albahaca dulce, los
tratamientos DBD redujeron \emph{Fusarium oxysporum} pero colonias
viables persistieron en los huecos de las semillas \textbf{{[}114{]}}.

\subsubsection{3.3.3. Respuestas fisiológicas y
moleculares}\label{respuestas-fisioluxf3gicas-y-moleculares}

Los RONS generados por plasma promueven el catabolismo del ácido
abscísico (ABA), mejoran la señalización del ácido giberélico (GA),
interactúan con el óxido nítrico (NO) e inducen la expresión de genes
debilitantes de la pared celular \textbf{{[}116, 117{]}}.

A nivel fisiológico, uno de los resultados más consistentes es la
sobre-regulación de las defensas antioxidantes. En colza expuesta a
sequía, el CAP mejoró la germinación y aumentó las actividades de SOD y
CAT mientras reducía la peroxidación lipídica (MDA) (Figura 12)
\textbf{{[}118{]}}.

\begin{quote}
\includegraphics{./images/image12.jpeg}
\end{quote}

\begin{quote}
\textbf{\emph{Figura 12. Efecto del estrés por sequía y el tratamiento
con plasma frío sobre 4 biomarcadores en plántulas de colza: (a)
actividad SOD, (b) actividad CAT, (c) contenidos de MDA, (d) contenidos
de azúcares solubles. WW: bien regado; DS: estrés por sequía. Adaptado
de {[}118{]}.}}
\end{quote}

En arroz bajo estrés por frío, el DBD de argón aceleró la germinación y
activó cambios transcriptómicos incluyendo la supresión de NCED3 y la
inducción de genes de tolerancia al estrés \textbf{{[}119{]}}.
Tendencias similares se reportaron en soja \textbf{{[}120{]}}.

\subsection{3.4. Efectos del plasma frío en plántulas y
plantas}\label{efectos-del-plasma-fruxedo-en-pluxe1ntulas-y-plantas}

El CAP puede estimular el crecimiento y la tolerancia al estrés en
múltiples especies, aunque los resultados permanecen fuertemente
dependientes de la dosis y el contexto. El PAW promovió el crecimiento
de plántulas de tomate y la calidad nutricional, mientras que la
albahaca expuesta a DBD mostró clorofilas elevadas y actividad
antioxidante \textbf{{[}123, 124{]}}. Los beneficios a largo plazo del
riego con PAW son visibles en tomate y pimiento, donde las plántulas
permanecieron más grandes y vigorosas incluso 60 días después de la
siembra (Figuras 13a--b) \textbf{{[}125{]}}.

En tomate, los tratamientos moderados de PAW (15--30 min) maximizaron el
desarrollo de plántulas a los 28 días, mientras que exposiciones más
largas se volvieron menos efectivas (Figura 13c) \textbf{{[}123{]}}. Las
plántulas de lenteja regadas con agua activada por plasma mostraron un
crecimiento más fuerte comparado con controles a los 18 días (Figura
13d) \textbf{{[}126{]}}.

\begin{quote}
\includegraphics{./images/image13.jpeg}
\end{quote}

\begin{quote}
\textbf{\emph{Figura 13. (a--b) Efecto a largo plazo del PAW sobre
plántulas de tomate (a) y pimiento (b), 60º día después de la siembra
{[}125{]}. (c) Efecto del tiempo de exposición al PAW sobre el
crecimiento de plántulas de tomate (28 días) {[}123{]}. (d) Plántulas de
lenteja regadas con agua del grifo (TW), agua del grifo activada por
plasma (PATW), agua desmineralizada (DM) o agua desmineralizada activada
por plasma (PADM), día 18 {[}126{]}.}}
\end{quote}

El plasma también afecta el rendimiento fotosintético (PSII). Las dosis
bajas mejoran los parámetros fotoquímicos, mientras que las dosis más
altas los inhiben \textbf{{[}127, 128{]}}. CAP/PAW a menudo sobre-regula
enzimas antioxidantes (SOD, CAT, POX) y modula el metabolismo secundario
\textbf{{[}124, 129{]}}. Las plántulas de tomate rociadas con PAW
mostraron hasta un 61\% de reducción en severidad de enfermedad contra
\emph{Xanthomonas} \textbf{{[}130{]}}.

La química redox-activa del PAW se ha vinculado a una mejor tolerancia a
la salinidad y la sequía \textbf{{[}128, 132{]}}. A nivel radicular, el
maíz mostró lignificación endodérmica acelerada bajo exposición al
arsénico, sugiriendo una ``fortificación'' anatómica inducida por plasma
\textbf{{[}133{]}}.

\begin{center}\rule{0.5\linewidth}{0.5pt}\end{center}

\section{4. Conclusiones y
perspectivas}\label{conclusiones-y-perspectivas}

Este artículo ha trazado una línea continua, aunque discontinua, desde
las ``electrizaciones'' del siglo XVIII hasta la agronomía de plasma
frío de hoy. El eje es la metrología. Donde Bertholon diseñó caminos
para un ``fluido eléctrico'' indivisible sin acceso a dosis ni agentes,
la agricultura plasma contemporánea define fuentes, cuantifica
transferencias interfaciales y mide la dosis bajo condiciones
controladas.

Reevaluada a la luz de esto, la contribución de Bertholon es
metodológica más que evidencial: proporcionó un vocabulario operativo y
la intuición de que regímenes eléctricos moderados podrían modular
benéficamente las funciones de las plantas. El trabajo moderno explica
cuándo y por qué esa intuición se mantiene.

La ruptura con la electrocultura histórica es triple: la naturaleza del
objeto eléctrico (de una esencia unitaria a un proceso multifísico), los
medios de medición y control (de señales cualitativas a dosimetría
químico-energética) y los estándares de validación (de pruebas
convincentes a diseños aleatorizados y replicados).

De cara al futuro, cuatro prioridades definen un agenda aplicada:

\begin{enumerate}
\def\labelenumi{\arabic{enumi}.}
\tightlist
\item
  \textbf{Anclar la práctica en dosificación rastreable} para ambas
  rutas CAP y PAW
\item
  \textbf{Escalar} con fuentes eco-diseñadas, alimentadas por aire y
  reactores PAW de baja energía
\item
  \textbf{Integrar procesos de plasma} dentro de itinerarios
  agroecológicos
\item
  \textbf{Garantizar una adopción segura y robusta} mediante límites de
  ozono/NOₓ y UV, protección contra arcos y ensayos multicéntricos
\end{enumerate}

Visto a lo largo de tres siglos, el motivo de una ``électricité
vivifiante'' perdura, pero su estatus está transformado. Ya no es una
idea reguladora basada en reservorios atmosféricos; es una caja de
herramientas mecanicista capaz de entregar sistemas de cultivo de menor
insumo, rastreables y potencialmente más resilientes.

\begin{center}\rule{0.5\linewidth}{0.5pt}\end{center}

\section{5. Referencias}\label{referencias}

{[}1{]} Teofrasto. \emph{De Lapidibus}. Oxford, At the Clarendon Press
(1965). {[}2{]} Plinio el Viejo. \emph{Historia Natural}. Libro 37.
Traducido por D.E. Eichholz (1962). {[}3{]} J. Francis, J. Dingley.
``Electroanaesthesia --- del pez torpedo al TENS.'' \emph{Anaesthesia}
70(1), 93--103 (2015). {[}4{]} A. Magnus, \emph{De mineralibus}, Libro
II. Clarendon Press: Oxford, 1967. {[}5{]} G. della Porta. \emph{Magia
naturalis} (1558/1589), Libro VII, cap. 16. {[}6{]} G. Cardano. \emph{De
subtilitate}, Libro VII. Lyon, 1663. {[}7{]} W. Gilbert. \emph{De
Magnete} (1600). {[}8{]} N. H. de V. Heathcote, ``El globo de azufre de
Guericke,'' \emph{Ann. Sci.} 6(3), 293--305 (1950). {[}9{]} F. Hauksbee.
\emph{Experiencias Físico-Mecánicas sobre Varios Temas} (1709). {[}10{]}
S. Gray. ``Una carta a Cromwell Mortimer.'' \emph{Philosophical
Transactions} 37, 1731--32. {[}11{]} C.-F. Du Fay. ``Electricidad
Vitrea/Resinosa.'' \emph{Nature} 144, 105 (1939). {[}12{]} J. L.
Heilbron. ``A propos de l'invention de la bouteille de Leyde.''
\emph{Revue d'histoire des sciences} 19(2), 133--142 (1966). {[}13{]} C.
Dorsman, C. A. Crommelin. ``La invención de la botella de Leyden.''
{[}14{]} J. Priestley. \emph{The History and Present State of
Electricity}. Londres, 1767. {[}15{]} E. Nairne. \emph{Experiencias
sobre electricidad}. Londres, 1779. {[}16{]} W. D. Hackmann. ``Las
investigaciones del Dr.~Martinus Van Marum.'' \emph{Medical History}
16(1), 11--26 (1972). {[}17{]} J.-A. Nollet. \emph{Cartas sobre la
electricidad}. París, 1753. {[}18{]} J.-A. Nollet. \emph{Ensayo sobre la
Electricidad de los Cuerpos}. París, 1765. {[}19{]} B. Pike.
\emph{Catálogo Ilustrado Descriptivo de Pike}, 2ª ed., vol.~1, 1856,
pp.~260, 269. {[}20{]} A. Ganot. \emph{Tratado elemental de física
experimental}, 9ª ed.~(1860), pp.~567, 691. {[}21{]} B. Franklin.
\emph{Experiencias y observaciones sobre la electricidad}. Londres,
1751. {[}22{]} E. P. Krider. ``Benjamin Franklin y los pararrayos.''
\emph{Physics Today} 59(1), 42--48 (2006). {[}23{]} T. Lane.
``Descripción de un electrómetro inventado por el Sr.~Lane.''
\emph{Phil. Trans. R. Soc.} 57, 451--460 (1767). {[}24{]} J.-A. Sigaud
De La Fond. \emph{Resumen histórico y experimental de los fenómenos
eléctricos}. París, 1785. {[}25{]} \emph{Philosophical Transactions},
vol.~LXII, 1772, p.~359. {[}26{]} J. Lacki. ``El Turista Físico.
Ginebra.'' \emph{Physics in Perspective} 9(2), 231--252 (2007). {[}27{]}
C. A. Coulomb. \emph{Memorias sobre la electricidad y el magnetismo}.
París, 1785. {[}28{]} J. Priestley. \emph{The History and Present State
of Electricity}, 2ª ed.~Londres, 1769. {[}29{]} T. Cavallo. \emph{Un
tratado completo de electricidad}, 3ª ed.~Londres, 1786. {[}30{]} G.
Beccaria. \emph{Dell'elettricismo artificiale e naturale}. 1753.
{[}31{]} Bertholon. \emph{De l'électricité des végétaux}. París, 1783.
{[}32{]} L. Fabbrizzi. ``Extraño Caso del Señor Volta y el Señor
Nicholson.'' \emph{Angewandte} 58(18), 5810--5822 (2019). {[}33{]} J. C.
Bose. \emph{Electrofisiología Comparada}. Nueva York, 1907. {[}34{]} W.
Crookes. ``Sobre la materia radiante.'' \emph{American Journal of
Science} s3-18(106) (1879). {[}35{]} W. von Siemens. ``Sobre la
inducción electrostática.'' \emph{Annalen der Physik} 178(9), 66--122
(1857). {[}36{]} M. Faraday. ``Investigaciones experimentales en
electricidad --- quinta serie.'' \emph{Philosophical Transactions} 123
(1833). {[}37{]} J. Burdon-Sanderson. ``Nota sobre los Fenómenos
Eléctricos que acompañan la Irritación de la Hoja de Dionaea
muscipula.'' \emph{Proc. Royal Society} (1873). {[}38{]} F. Paschen.
``Sobre la diferencia de potencial requerida para la chispa en aire.''
\emph{Annalen der Physik} (1889). {[}39{]} J. S. Townsend. ``La
conductividad producida en gases.'' \emph{Nature} 62, 340--341 (1900).
{[}40{]} J. S. Townsend. \emph{La Teoría de Ionización de Gases por
Colisión}. Londres, 1901. {[}41{]} A. W. Hofmann. ``Dr.~Heinrich
Geissler.'' \emph{Ber. Dtsch. Chem. Ges.} 12, 147--148 (1879). {[}42{]}
I. Langmuir. ``Oscilaciones en Gases Ionizados.'' \emph{Proc. Natl.
Acad. Sci.} (1928). {[}43{]} C. von Sonntag, U. von Gunten.
\emph{Química del ozono en tratamiento de agua y aguas residuales}. Iwa
Publishing, 2012. {[}44{]} C. Gottschalk, J. A. Libra, A. Saupe.
\emph{Ozonación de agua y aguas residuales}. Wiley-VCH, 2010. {[}45{]}
Kr. Birkeland. \emph{Trans. Faraday Soc.} 2, 98--116 (1906). {[}46{]} A.
M. W. Downing, T. P. Blunt. ``Investigaciones sobre el efecto de la luz
sobre Bacterias.'' \emph{Proc. Royal Society of London} 26 (1878).
{[}47{]} M. Luckiesh, L. Holladay. \emph{Aplicaciones de Energía
Germicida, Eritemal e Infrarroja}. Nueva York, 1946. {[}48{]}
Rentschler, H. (1942) ``Acción Bactericida de la Radiación
Ultravioleta.'' {[}49{]} S. Lemström. \emph{Electricidad en Agricultura
y Horticultura}. Londres, 1904. {[}50{]} J. Christofleau.
\emph{Electroculture --- Fertilización de tierras por la electricidad
atmosférica}. 1925. {[}51{]} M. S. Kyi, J. Holton, G. L. Ridgway.
``Evaluación de la eficacia de un sistema de esterilización por plasma
gaseoso de H₂O₂ a baja temperatura.'' \emph{J. Hospital Infection}
31(4), 275--284 (1995). {[}52{]} W. A. Rutala et al.~``Evaluación
comparativa de la actividad esporicida de nuevas tecnologías de
esterilización a baja temperatura.'' \emph{Am. J. Infection Control}
26(4), 393--398 (1998). {[}53{]} M. Laroussi. ``Esterilización de
materia contaminada con un plasma a presión atmosférica.'' \emph{IEEE
Trans. Plasma Science} 24(3), 1188--1191 (1996). {[}54{]} I. Adamovich
et al.~``La Hoja de Ruta del Plasma 2017.'' \emph{J. Physics D: Applied
Physics} 50(32), 323001 (2017). {[}55{]} P. Ranieri et al.~``Agricultura
plasma: Revisión.'' \emph{Plasma Processes \& Polymers} 18(1), 2000162
(2021). {[}56{]} D. B. Graves. ``El papel emergente de las especies
reactivas de oxígeno y nitrógeno.'' \emph{J. Physics D: Applied Physics}
45(26), 263001 (2012). {[}57{]} D. Yan et al.~``Mejorando la Germinación
de Semillas con Plasma Atmosférico Frío.'' \emph{Plasma} 5(1), 98--110
(2022). {[}58{]} M. A. Benabderrahm et al.~``El tratamiento con plasma
frío impulsa la germinación de cebada.'' \emph{J. Cereal Science} 116,
103852 (2024). {[}59{]} K. S. Wong et al.~``Agua Activada por Plasma:
Propiedades Fisicoquímicas.'' \emph{Processes} 11(7), 2213 (2023).
{[}60{]} E. Cortese et al.~``PAW Provoca Elevaciones Rápidas de Ca²⁺ en
Arabidopsis.'' \emph{Plants} 10(11), 2516 (2021). {[}61{]} S. M. E.
Sultan et al.~``El plasma atmosférico frío mejora plántulas de tomate.''
\emph{BMC Plant Biology} 24, 420 (2024). {[}62{]} Y. Zambon et al.~``El
agua activada por plasma provoca respuestas de defensa en plantas.''
\emph{Scientific Reports} 10, 19211 (2020). {[}63{]} Y. Gao et
al.~``Revisión sobre la formación de agua activada por plasma frío.''
\emph{Food Research International} 157, 111246 (2022). {[}64{]} J.
Šimečková et al.~``Influencia del PAW sobre propiedades del suelo.''
\emph{Water} 12(9), 2357 (2020). {[}65{]} B. A. Niemira.
``Descontaminación de Alimentos con Plasma Frío.'' \emph{Ann. Rev.~Food
Sci. Tech.} 3, 125--142 (2012). {[}66{]} N. N. Misra et
al.~``Tratamiento con plasma frío a presión atmosférica en empaque de
fresas.'' \emph{J. Food Engineering} 125, 131--138 (2014). {[}67{]} M.
Bayati et al.~``Cambios químicos y físicos inducidos por tratamiento con
plasma frío.'' \emph{Compr Rev Food Sci Food Saf} 23(4), e13376 (2024).
{[}68{]} M. François, F. Fourmaux. ``Nota sobre el Abbé Pierre Bertholon
de Saint Lazare.'' CTHS. {[}69{]} L. Dulieu. ``El movimiento científico
montpellierés en el siglo XVIII.'' \emph{Revue d'histoire des sciences}
11(3), 227--249 (1958). {[}70{]} M. l'Abbé Bertholon. \emph{De
l'électricité du corps humain}. París, 1780. {[}71{]} M. l'Abbé
Bertholon, M. Goyer. \emph{La nature considérée sous ses différents
aspects}. París, 1787. {[}72{]} J. Priestley. ``Una cuenta de un nuevo
electrómetro, ideado por el Sr.~William Henly.'' \emph{Philosophical
Transactions} 62 (1772). {[}73{]} A. Guillemin. \emph{Electricidad y
Magnetismo}. Londres, 1891. {[}74{]} C. A. Coulomb. \emph{Memorias sobre
la electricidad y el magnetismo}. Segunda memoria, pp.~577--611 (1785).
{[}75{]} A. Martinez. ``Replicación del Experimento de la Balanza de
Torsión de Coulomb.'' \emph{Arch. Hist. Exact Sci.} 60, 517--563 (2006).
{[}76{]} M. B. Rubin. ``La historia del ozono.'' \emph{Bull. Hist.
Chem.} 26(1), 40--56 (2001). {[}77{]} D. Tarasick et al.~``Informe de
Evaluación de Ozono Troposférico.'' \emph{Elementa} 7:39 (2019).
{[}78{]} J. Barth. ``Johann Wilhelm Ritter y el descubrimiento de la
radiación UV.'' \emph{Hautarzt} 38(5), 301--303 (1987). {[}79{]} E.
Kuffel et al.~\emph{Ingeniería de alto voltaje: fundamentos}.
Butterworth-Heinemann, 2000. {[}80{]} R. A. Fisher. \emph{El diseño de
experimentos}. Hafner Press, 1971. {[}81{]} I. Chalmers et al.~``Por qué
el ensayo MRC de 1948 de estreptomicina usó asignación basada en números
aleatorios.'' \emph{J. Royal Society of Medicine} 104(9) (2011).
{[}82{]} P. Bourke et al.~``El Potencial del Plasma Frío para una
Producción Alimentaria Segura y Sostenible.'' \emph{Trends in
Biotechnology} 36(6), 615--626 (2018). {[}83{]} F. F. Chen.
\emph{Introducción a la Física del Plasma y la Fusión Controlada}.
Springer, 2019. {[}84{]} Y. P. Raizer. \emph{Física de Descargas de
Gas}. Springer, 1991. {[}85{]} M. A. Lieberman, A. J. Lichtenberg.
\emph{Principios de Descargas de Plasma y Procesamiento de Materiales}.
Wiley, 2005. {[}86{]} L. C. Pitchford et al.~``LXCat: una Plataforma de
Acceso Abierto.'' \emph{Plasma Processes \& Polymers} 14(1--2) (2016).
{[}87{]} F. Judée et al.~``Activación de agua del grifo usando DBD.''
\emph{Wat. Res.} 133, 47--59 (2018). {[}88{]} P. Bruggeman, C. Leys.
``Plasmas no térmicos en y en contacto con líquidos.'' \emph{J. Physics
D} 42(5), 053001 (2009). {[}89{]} P. J. Bruggeman et al.~``Interacciones
plasma-líquido: Una revisión y hoja de ruta.'' \emph{Plasma Sources Sci.
Technol.} 25(5), 053002 (2016). {[}90{]} U. Kogelschatz. ``Descargas de
Barrera Dieléctrica.'' \emph{Plasma Chemistry and Plasma Processing} 23,
1--46 (2003). {[}91{]} J. Winter et al.~``Chorros de plasma a presión
atmosférica: una visión general.'' \emph{Plasma Sources Sci. Technol.}
24(6), 064001 (2015). {[}92{]} T. C. Manley. ``Las Características
Eléctricas de la Descarga del Ozonizador.'' \emph{Trans. Electrochemical
Society} 84(1) (1943). {[}93{]} H. Decauchy, T. Dufour. ``Mecanismos de
transmisión y reflexión múltiple de streamers guiados.'' \emph{Plasma
Sources Sci. Technol.} 31(11) (2022). {[}94{]} J. H. Kim et
al.~``Mediciones ópticas de temperaturas de gas en plasmas fríos RF a
presión atmosférica.'' \emph{Surface and Coatings Technology} 171,
211--215 (2003). {[}95{]} N. Masoud et al.~``Mediciones de Temperaturas
Rotacionales y Vibracionales en un C-DBD.'' \emph{Contributions to
Plasma Physics} 45(1), 32--39 (2005). {[}96{]} A. Herrmann et
al.~``Diseño Mejorado de Sonda Catalítica para Mapeo de Densidad de
Radicales.'' \emph{J. Phys. Chem. A} 128(46), 10080--10086 (2024).
{[}97{]} R. R. Brandenburg. ``Descargas de barrera dieléctrica: progreso
en fuentes de plasma.'' \emph{Plasma Sources Sci. Technol.} 26(5),
053001 (2017). {[}98{]} L. L. Alves et al.~``Fundamentos de estándares
de plasma.'' \emph{Plasma Sources Sci. Technol.} 32(2) (2023). {[}99{]}
J. R. Bolton, K. G. Linden. ``Estandarización de Métodos para
Determinación de Fluencia.'' \emph{J. Environmental Engineering} 129(3).
{[}100{]} D. Tsikas. ``Análisis de nitrito y nitrato en fluidos
biológicos.'' \emph{J. Chromatography B} 851, 51--70 (2007). {[}101{]}
J.-Y. Han et al.~``Efectos del agua activada por burbuja de plasma sobre
patógenos transmitidos por alimentos en tomates.'' \emph{Food Control}
144, 109381 (2023). {[}102{]} M. M. Rahman et al.~``Mecanismos y
Señalización Asociados con la Mejora del Crecimiento en Trigo Mediante
Plasma LPDBD.'' \emph{Scientific Reports} 8, 10498 (2018). {[}103{]} N.
Dawood. ``Efectos del tratamiento con plasma de aire sobre la absorción
de agua de semillas de trigo y cebada.'' \emph{J. Taibah University For
Science} 15(1), 1094--1100 (2021). {[}104{]} P. Starič et
al.~``Respuesta de Dos Variedades de Trigo a Plasma de Oxígeno de Brillo
y Afterglow.'' \emph{Plants} 10(8), 1728 (2021). {[}105{]} P. Starič et
al.~``La Influencia del Plasma Frío de Brillo y Afterglow sobre Semillas
de Trigo.'' \emph{Int. J. Mol. Sci.} 23(13), 7369 (2022). {[}106{]} S.
Karmakar et al.~``Impacto del plasma LFGD sobre superficie de semillas,
germinación, crecimiento de maíz.'' \emph{Heliyon} 7, e06458 (2021).
{[}107{]} R. Molina et al.~``Análisis fisicoquímico de superficie de
semillas de trigo tratadas con DBD.'' \emph{Plasma Processes \&
Polymers} 18(1), 2000086 (2021). {[}108{]} N. Recek et al.~``Germinación
de Semillas de Phaseolus vulgaris L. tras Plasma RF.'' \emph{Int. J.
Mol. Sci.} 22(13), 6672 (2021). {[}109{]} E. Bormashenko et al.~``El
Tratamiento con Plasma RF Frío Modifica la Humectabilidad y Velocidad de
Germinación de Semillas.'' \emph{Scientific Reports} 2, 741 (2012).
{[}110{]} N. Khamsen et al.~``Esterilización y Mejora de Germinación de
Semillas de Arroz Mediante Plasma de Descarga Híbrida Atmosférica.''
\emph{ACS Appl. Mater. Interfaces} 8(30), 19268--19275 (2016). {[}111{]}
J. Mravlje et al.~``Desarrollo de Tecnologías de Plasma Frío para
Descontaminación Superficial de Patógenos Fúngicos en Semillas.''
\emph{J. Fungi} 7(8), 650 (2021). {[}112{]} A. Mitra et
al.~``Inactivación de Microorganismos Superficiales por Plasma
Atmosférico Frío.'' \emph{Food and Bioprocess Technology} 7, 645--653
(2014). {[}113{]} P. Basaran et al.~``Eliminación de Aspergillus
parasiticus de superficie de nueces con LPCP.'' \emph{Food Microbiology}
25, 626--632 (2008). {[}114{]} K. Homa et al.~``Estrategias de
Tratamiento con Plasma Frío para Fusarium oxysporum en Albahaca Dulce.''
\emph{HortScience} 56(1), 42--51 (2020). {[}115{]} J. Mravlje et
al.~``El Plasma Frío Afecta la Germinación y la Estructura de la
Comunidad Fúngica de Semillas de Trigo Sarraceno.'' \emph{Plants} 10(5),
851 (2021). {[}116{]} G. Grainge et al.~``Mecanismos moleculares de
liberación de dormancia de semillas por agua activada con plasma
gaseoso.'' \emph{J. Experimental Botany} 73(12), 4065--4078 (2022).
{[}117{]} G. Grainge et al.~``Impacto del Agua Activada con Plasma
Gaseoso sobre Mecanismos de Dormancia Fotodependiente en Semillas de
Nicotiana tabacum.'' \emph{Int. J. Mol. Sci.} 23(12), 6709 (2022).
{[}118{]} L. Ling et al.~``El tratamiento con plasma frío mejora la
germinación de semillas de colza bajo estrés por sequía.''
\emph{Scientific Reports} 5, 13033 (2015). {[}119{]} J.-Y. Bian et
al.~``El plasma no térmico mejora la germinación de semillas de arroz
bajo estrés por baja temperatura.'' \emph{Applied Biological Chemistry}
67(2) (2024). {[}120{]} K. Sayahi et al.~``Evaluando el impacto del
Plasma Frío sobre la soja.'' \emph{BMC Biotechnology} 24, 93 (2024).
{[}121{]} K. Kučerová et al.~``Efectos del agua activada por plasma
sobre el trigo.'' \emph{Plasma Processes \& Polymers} 16(3), 1800131
(2019). {[}122{]} B. Šerá et al.~``Tratamiento con Plasma de Arco
Deslizante de Granos de Maíz.'' \emph{Agronomy} 11(10), 2066 (2021).
{[}123{]} B. Adhikari et al.~``Riego con Agua Activada por Plasma
Atmosférico Frío en Plántulas de Tomate.'' \emph{Scientific Reports} 9,
16080 (2019). {[}124{]} F. M. Abarghuei et al.~``Una Aplicación de CAP
para Mejorar Rasgos de Albahaca.'' \emph{Plants} 10(10), 2088 (2021).
{[}125{]} L. Sivachandiran, A. Khacef. ``Germinación mejorada de
semillas y crecimiento de plantas por plasma de aire frío a presión
atmosférica.'' \emph{RSC Advances} 7, 1822 (2017). {[}126{]} S. Zhang,
A. Rousseau, T. Dufour. ``Promoviendo la germinación de lenteja con agua
activada por plasma.'' \emph{RSC Advances} 7, 31244 (2017). {[}127{]} J.
Perner et al.~``El tratamiento con plasma frío influye en los parámetros
fisiológicos del mijo.'' \emph{Photosynthetica} 62(1), 126--137 (2024).
{[}128{]} F. Bussmann et al.~``Efectos a Largo Plazo del Agua Tratada
con CAP sobre Hordeum vulgare.'' \emph{J Plant Growth Regul} 42,
3274--3290 (2023). {[}129{]} K. Panngom et al.~``El Tratamiento con
Plasma No Térmico Reduce la Viabilidad Fúngica.'' \emph{PLoS ONE} 9(6),
e99300 (2014). {[}130{]} S. M. Perez et al.~``El agua activada por
plasma como inductor de resistencia contra mancha bacteriana del
tomate.'' \emph{PLoS ONE} 14(5), e0217788 (2019). {[}131{]} P. J. Savi
et al.~``El riego con agua activada por plasma aumenta la mortalidad de
ácaros araña en tomate.'' \emph{Scientific Reports} 15, 22118 (2025).
{[}132{]} M. Veerana et al.~``El plasma no térmico mejora el crecimiento
y la tolerancia a la salinidad de bok choy.'' \emph{Frontiers in Plant
Science} 15, 1445791 (2024). {[}133{]} Z. Lukacova et al.~``El Efecto
del PAW sobre Maíz bajo Estrés por Arsénico.'' \emph{Plants} 10(9), 1899
(2021).

\begin{center}\rule{0.5\linewidth}{0.5pt}\end{center}

\begin{quote}
📄 \textbf{Paper original:}
\href{https://arxiv.org/abs/2505.06968}{arXiv:2505.06968} 📥
\textbf{Descargar PDF:}
\href{https://arxiv.org/pdf/2505.06968}{arxiv.org/pdf/2505.06968} 🔬
\textbf{DOI:}
\href{https://doi.org/10.5802/crmeca.331}{10.5802/crmeca.331}
\end{quote}

\end{document}
